\subsection*{Ex. 1.4 (5 pts) --- ``Segmentation fault (core dumped)''}
\label{p1:sec:segfault}

\textbf{What the message means.} The program touched memory it was not entitled to touch.
The MMU checks every address as it is translated, and when the address is either not
mapped at all or is mapped with the wrong permission (writing a read-only page, executing
a non-executable one), it raises a hardware fault. The kernel finds no legitimate reason
for the fault, converts it into the signal \texttt{SIGSEGV}, and delivers it to the
process; the default action for \texttt{SIGSEGV} is to kill the process and write out a
core dump. The line you see is the shell afterwards reporting how its child died.

\textbf{Possible reasons.}
\begin{itemize}[nosep,leftmargin=1.4em]
  \item Dereferencing a null pointer (typically an unchecked return from
        \texttt{malloc} or \texttt{fopen}).
  \item Dereferencing an uninitialized or otherwise wild pointer.
  \item Use-after-free, or following a dangling pointer to a stack frame that has
        already returned.
  \item Running past the end of an array or buffer, far enough to leave the mapping.
  \item Writing through a \texttt{char *} that points at a string literal, which lives on
        a read-only page.
  \item Unbounded recursion overflowing the stack into the guard page below it.
  \item Corrupting the allocator's metadata (double free, overrun of a heap block), so
        that a later allocation returns nonsense.
\end{itemize}

\textbf{What ``segmentation'' means here.} The word is historical. On machines with
segmented memory the address space was divided into segments, each with a base, a limit
and access rights, and the hardware trapped any reference that fell outside the bounds or
rights of the current segment --- a \emph{segmentation violation}. Mainstream systems are
paged rather than segmented now, but the name stuck, and it means what it always meant:
an access outside a valid, permitted region of the address space.

\textbf{What ``core'' means here.} ``Core'' is magnetic-core memory, the ferrite-ring
main memory of machines of the 1950s and 60s; ``core'' became a synonym for main memory
and outlived the technology. A core dump is therefore an image of the process's memory
--- its segments as listed in Ex.~1.2 --- together with its register state, captured at
the instant it faulted and written to a file so that a debugger can be pointed at it
afterwards and reconstruct where and why the program died.
